% =============================================================================
% CAPÍTULO 2 — Marco Teórico y Estado del Arte
% =============================================================================

\chapter{Marco Teórico y Estado del Arte}

\section{Mantenimiento Industrial, conceptos y evolución}

\section{Evolución histórica de las estrategias de mantenimiento}

\section{Conceptos fundamentales}
\subsection{Falla}
\subsection{Anomalía operativa}
\subsection{Tiempo de vida útil (RUL)}

\section{Series temporales y preprocesamiento para datos de sensores}
\subsection{Características de los datos de sensores industriales}

\section{Etapas de preprocesamiento de series temporales}
\subsection{Limpieza de datos}
\subsection{Normalización y estandarización}
\subsection{Validación respecto al orden temporal}

\section{Consideraciones específicas para estimar el RUL}

\section{Tareas de mantenimiento predictivo y consideraciones}
\subsection{Detección de anomalías}
\subsection{Tipos de anomalías en series de tiempo}
\subsection{Enfoques para la detección de anomalías}
\subsection{Dificultades para la detección de anomalías}
\subsection{Estimación de RUL}

\section{Algoritmos de referencia}
\subsection{Random Forest}
\subsection{XGBoost}
\subsection{LSTM}
\subsection{Isolation Forest}

\section{Predictor Gamma}
\subsection{Características principales}
\subsection{Operadores fundamentales}
\subsection{Codificación JM}
\subsection{Operador de similitud}
\subsection{Notación}
\subsection{Algoritmo}

\section{Conjuntos de datos}
\subsection{Dificultades para obtener datos reales}
\subsection{C-MAPSS}

\section{Medidas de desempeño}
\subsection{Estimación del RUL}
\subsection{Detección de anomalías}

\section{Análisis estadístico}
\subsection{Prueba de Friedman}
\subsection{Test post-hoc de Nemenyi}
\subsection{Corrección de Bonferroni}

\section{Estado del arte}
